% =====================================================================
%  CHAPTER 11: 学校健康教育
% =====================================================================
\chapter{学校健康教育}
\setcounter{chapter}{11}
\chapterintro{
掌握学校健康教育的概念和基本内容；掌握学校健康教育实施的主要形式和评价方法；掌握健康教育在学校卫生工作中的应用；熟悉专题健康教育（营养教育、青春期性教育、心理健康教育、安全教育等）；了解健康促进学校（HPS）的创建。
}

% ----- 11.1 概述 -----
\section{学校健康教育概述}

\begin{defbox}
\textbf{学校健康教育（School Health Education）}：在学校环境中，通过有目的、有计划、有组织的教育教学活动，向学生传授健康知识、培养健康行为习惯、树立健康价值观念，以提高学生健康素养，促进其身心全面健康发展的系统教育过程。

\textbf{学校健康教育的意义：}
\begin{enumerate}[leftmargin=*]
    \item 学生时期是行为和生活方式形成的关键期
    \item 学校是开展健康教育最理想的场所（覆盖面广、持续性强、影响深远）
    \item 学生可将健康知识带回家，产生"小手拉大手"的家庭辐射效应
    \item 预防胜于治疗——通过教育预防健康问题的发生
\end{enumerate}
\end{defbox}

\subsection{学校健康教育的基本内容}

\begin{keybox}
\textbf{学校健康教育的六大主题内容：}

\begin{enumerate}[leftmargin=*]
    \item \imp{健康行为与生活方式}
    \begin{itemize}
        \item 合理膳食与营养、规律运动、充足睡眠
        \item 个人卫生习惯（口腔卫生、用眼卫生、手卫生）
        \item 拒绝吸烟、饮酒和物质滥用
    \end{itemize}

    \item \imp{疾病预防}
    \begin{itemize}
        \item 常见传染病预防（流感、手足口病等）
        \item 常见病预防（近视、龋齿、肥胖等）
        \item 慢性病早期预防
    \end{itemize}

    \item \imp{心理健康}
    \begin{itemize}
        \item 情绪管理、压力应对
        \item 人际交往技能
        \item 自我认知与自尊
    \end{itemize}

    \item \imp{生长发育与青春期保健}
    \begin{itemize}
        \item 生长发育知识
        \item 青春期生理心理变化
        \item 生殖健康与性教育
    \end{itemize}

    \item \imp{安全应急与避险}
    \begin{itemize}
        \item 交通安全、消防安全、防溺水
        \item 自然灾害应对
        \item 校园欺凌预防
    \end{itemize}

    \item \imp{健康素养与健康责任}
    \begin{itemize}
        \item 培养健康意识和健康责任
        \item 提高健康信息获取和判别能力
    \end{itemize}
\end{enumerate}
\end{keybox}

% ----- 11.2 专题教育 -----
\section{专题健康教育}

\subsection{营养教育}

\begin{defbox}
\textbf{营养教育目标：}使学生了解营养素的功能和食物来源，掌握合理膳食的基本原则，养成健康的饮食行为和习惯。

\textbf{核心内容：}食物多样、谷物为主；多吃蔬菜水果和薯类；每天吃奶类大豆或其制品；常吃适量鱼禽蛋瘦肉；减少烹调油用量；吃新鲜卫生食物；三餐定时定量、不偏食挑食。
\end{defbox}

\subsection{青春期性教育}

\begin{defbox}
\textbf{青春期性教育目标：}帮助青少年正确认识和对待青春期性发育，获得正确的性知识和生殖健康信息，形成健康的性观念和负责任的行为。

\textbf{核心内容：}青春期生理变化与个人卫生、性发育的正常过程、月经和遗精的卫生保健、性传播疾病预防、性心理发展、性别平等与尊重。
\end{defbox}

\subsection{心理健康教育}

\begin{defbox}
\textbf{心理健康教育目标：}提高学生心理素质，培养积极乐观、健康向上的心理品质，促进学生身心和谐可持续发展。

\textbf{核心内容：}认识和管理情绪、建立良好人际关系、提高抗挫折能力、自我认知与接纳、求助意识和能力。
\end{defbox}

\subsection{安全教育}

\begin{defbox}
\textbf{安全教育目标：}提高学生的安全防范意识和自我保护能力，减少意外伤害的发生。

\textbf{核心内容：}交通安全、防火防电防溺水、防校园欺凌和暴力、急救知识（心肺复苏等）、自然灾害逃生、网络安全与防沉迷。
\end{defbox}

% ----- 11.3 实施与评价 -----
\section{学校健康教育的实施}

\subsection{实施的主要形式}

\begin{keybox}
\textbf{学校健康教育实施的多种形式：}

\begin{enumerate}[leftmargin=*]
    \item \imp{健康教育课程}——独立设课或融入其他学科教学（如体育与健康课），是最系统和最基础的实施形式
    \item \imp{健康主题活动}——健康讲座、主题班会、健康知识竞赛、健康手抄报等
    \item \imp{同伴教育}——由同龄人向同龄人传递健康信息，在青春期性教育和心理健康教育中效果显著
    \item \imp{校园健康环境营造}——创建有利于健康行为形成的氛围（如健康食堂、无烟校园）
    \item \imp{新媒体与信息技术}——利用微信、短视频等新媒体传播健康知识
    \item \imp{家校合作}——通过家长会、亲子活动等向家长传播健康知识
\end{enumerate}
\end{keybox}

\subsection{评价方法}

\begin{defbox}
\textbf{学校健康教育的评价类型：}

\begin{enumerate}[leftmargin=*]
    \item \imp{过程评价（Process Evaluation）}
    \begin{itemize}
        \item 评价教育活动的实施过程和质量
        \item 指标：健康教育课时数、覆盖率、参与率、教材使用率
    \end{itemize}

    \item \imp{效果评价（Outcome Evaluation）}
    \begin{itemize}
        \item 评价教育活动的直接效果
        \item 指标：健康知识知晓率、健康态度正确率、健康行为形成率的变化
    \end{itemize}

    \item \imp{影响评价（Impact Evaluation）}
    \begin{itemize}
        \item 评价教育活动的长期健康影响
        \item 指标：常见病发病率、超重肥胖率、体能达标率的变化
    \end{itemize}
\end{enumerate}

\textbf{常用评价方法：}
\begin{itemize}[leftmargin=*]
    \item 问卷调查（KAP调查——知识、态度、行为）
    \item 焦点小组访谈
    \item 行为观察
    \item 健康监测数据的分析
    \item 前后对照设计
\end{itemize}
\end{defbox}

% ----- 11.4 健康促进学校 -----
\section{健康促进学校}

\begin{defbox}
\textbf{健康促进学校（Health Promoting School, HPS）}：WHO倡导的综合性学校健康促进模式，通过学校所有成员（学生、教职员工、家长和社区）的共同努力，提供一个完整、持续的经验来促进学生的健康。

\textbf{HPS的六大核心领域（WHO）：}
\begin{enumerate}[leftmargin=*]
    \item \imp{学校健康政策}——制定有利于健康的学校规章制度（如禁烟政策、食品安全政策）
    \item \imp{学校物质环境}——安全的校舍、清洁的饮用水、卫生设施、运动场所
    \item \imp{学校社会环境}——尊重包容的校园文化、良好的师生关系
    \item \imp{学校健康技能教育}——健康素养的培养
    \item \imp{学校与社区的联系}——与家长和社区建立健康合作关系
    \item \imp{学校卫生服务}——健康检查、常见病管理、心理咨询
\end{enumerate}
\end{defbox}

% ----- 复习题 -----
\section{复习题}

\begin{enumerate}[leftmargin=*, itemsep=8pt]
    \item \textbf{【名词解释】}学校健康教育；健康促进学校（HPS）
    \item \textbf{【名词解释】}KAP调查；同伴教育
    \item \textbf{【简答题】}简述学校健康教育的基本内容（六大主题）。
    \item \textbf{【简答题】}简述学校健康教育的实施形式。
    \item \textbf{【简答题】}简述学校健康教育的评价类型及主要方法。
    \item \textbf{【简答题】}简述健康促进学校（HPS）的六大核心领域。
\end{enumerate}

\subsection*{参考答案要点}

\begin{enumerate}[leftmargin=*, itemsep=5pt]
    \item \textbf{学校健康教育}：在学校环境中，通过有目的、有计划的教育教学活动，传授健康知识、培养健康行为、树立健康价值观的系统教育过程。\textbf{HPS}：WHO倡导的综合性学校健康促进模式，通过学校全体成员共同努力提供完整持续的健康促进经验。

    \item \textbf{KAP调查}：知识（Knowledge）、态度（Attitude）、行为（Practice）调查，是健康教育评价的常用方法。\textbf{同伴教育}：由同龄人向同龄人传递健康信息的教育方法，在青春期教育中效果显著。

    \item 六大主题：(1)健康行为与生活方式（营养、运动、卫生、拒烟酒）；(2)疾病预防（传染病、常见病、慢性病预防）；(3)心理健康（情绪管理、人际交往、自尊）；(4)生长发育与青春期保健（发育知识、性教育）；(5)安全应急与避险（交通、消防、防溺水、校园欺凌）；(6)健康素养与健康责任。

    \item 六种主要形式：(1)健康教育课程（最系统基础）；(2)健康主题活动（讲座、班会、竞赛）；(3)同伴教育；(4)校园健康环境营造；(5)新媒体与信息技术；(6)家校合作。

    \item 三种类型：过程评价（实施过程和质量）、效果评价（知识-态度-行为变化）、影响评价（长期健康指标变化）。常用方法：KAP问卷调查、焦点小组访谈、行为观察、健康监测数据分析、前后对照设计。

    \item 六大领域：(1)学校健康政策；(2)学校物质环境；(3)学校社会环境；(4)学校健康技能教育；(5)学校与社区的联系；(6)学校卫生服务。
\end{enumerate}

\newpage
